% source: An algebraic approach to the Riemann-Roch theorem and the arithmetic theory of function fields (Athanasios Papaioannou), http://www.fen.bilkent.edu.tr/~franz/mat/Riemann-Roch.pdf
% pdf_page: 0014
% retrieved: 2026-07-26
% transcribed_by: page-transcriber (AI vision, no OCR)

\documentclass[11pt]{article}
\usepackage{amsmath,amssymb,amsthm}

% Small-caps theorem style matching the source.
\newtheoremstyle{scplain}{\topsep}{\topsep}{\itshape}{0pt}{\scshape}{.}{ }{\thmname{#1} \thmnote{#3}}
\theoremstyle{scplain}
\newtheorem*{conjecture}{Conjecture}
\newtheorem*{theorem}{Theorem}

\DeclareMathOperator{\Supp}{Supp}

\begin{document}
\setlength{\emergencystretch}{2em}

This conjecture can be adapted suitably for number fields and has far-reaching
consequences, c.f. Lang [LA], Chapter IV, section 7, for a relevant discussion
and a number of references.

We shall now reformulate ABC so that its function field analog becomes apparent.
Before that, however, we will need a piece of ubiquitous terminology: we define
the \emph{height} of a rational number $r=m/n$ (where $(m,n)=1$) to be
$\operatorname{ht}(r)=\max\{\log(|m|),\log(|n|)\}$. Now we have

\begin{conjecture}[2.2 (ABC II)]
Given $\epsilon>0$, there is a constant $m_\epsilon$ such that for all rational
numbers $v$ and $u$ that satisfy the equation $v+u=1$ the following relation
holds:
\[
  \max\{\operatorname{ht}(v),\operatorname{ht}(u)\}
    \leq m_\epsilon+(1+\epsilon)\sum_{p\mid ABC}\log(p).
\]
\end{conjecture}

In the function field case, we will need a definition analogous to that of
height above; let $\mathbb{F}/\mathbb{K}$ be a function field, and let
$u \in \mathbb{F}^*$ be non-constant. Then, if $M$ is the maximal separable
extension of $\mathbb{K}(u)$ inside $\mathbb{F}$, we shall call the degree
$[M:\mathbb{K}(u)]$ of the extension $M/\mathbb{K}(u)$ the \emph{separable
degree of $u$}; it will be denoted by $\operatorname{ord}_s(u)$. We now have

\begin{theorem}[2.1 (ABC)]
Let $\mathbb{F}/\mathbb{K}$ be a function field. Then for
$v,u \in \mathbb{F}^*$ that satisfy the equation $v+u=1$ the following
inequality holds
\[
  \deg_s(v)=\deg_s(u) \leq 2g_{\mathbb{K}}
    + \sum_{P \in \Supp(A+B+C)}\deg_{\mathbb{K}}P,
\]
where $A,B$ are the zero divisors of $v$ and $u$ in $\mathbb{F}$ and $C$ is
their common pole divisor.
\end{theorem}

The ABC conjecture for function fields implies the analogs of all classical
number theoretic conjectures the classical ABC conjecture implies; in
particular, it implies Fermat's Last Theorem for function fields.

An important result that used in the proof of ABC is the analog of the Hurwitz
formula:

\begin{theorem}[2.2 (Hurwitz)]
Let $\mathbb{L}/\mathbb{K}$ be a finite separable extension of function fields.
Then,
\[
  2g_{\mathbb{L}}-2
    = [\mathbb{L}:\mathbb{K}](2g_{\mathbb{K}}-2)
      + \deg_{\mathbb{L}}D_{\mathbb{L}/\mathbb{K}};
\]
more generally,
\[
  2g_{\mathbb{L}}-2
    \geq [\mathbb{L}:\mathbb{K}](2g_{\mathbb{K}}-2)
      + \sum_{\mathfrak{p}}(e(\mathfrak{p}/P)-1)\deg_{\mathbb{L}}\mathfrak{p},
\]
where the sum is over all primes $\mathfrak{p}$ over $\mathbb{L}$ and equality
occurs if and only if all ramified primes are tamely ramified.
\end{theorem}

All the undefined terms appearing above have meaning analogous to their meaning
in the classical algebraic number theoretic context.

The above examples, however, appear meagre when compared to the other
achievements of this arithmetic theory for function fields; among its proved
theorems belong the Riemann Zeta Hypothesis (there is a natural, but advanced,
proof due to Weil [WE] and an elementary one by Bombieri -- this appears as an
Appendix to Rosen's book [RO]), the Artin Primitive Roots Conjecture, and the
Brumer-Stark Conjectures (there are two proofs, one by Tate and Deligne [TA] and
another by Hayes [HA]). As every other legitimate theory, it also has open
problems: the Langlands program has not been fully developed in the language of
function fields (despite important work in this direction by Drinfeld in the
1970's) and the Birch and Swinnerton-Dyer conjecture remains a mystery. The
methods one uses to formulate such conjectures and attain their proofs are not
much different than the ones used here, and it's probable that results unproved
in context other than that of a function field might be solvable in full
generality by the same methods used in the function field case \emph{mutatis
mutandis}.

\begin{center}14\end{center}

\end{document}
